%%
%% Copyright 2019-2020 Elsevier Ltd
%%
%% This file is part of the 'CAS Bundle'.
%% --------------------------------------
%%
%% It may be distributed under the conditions of the LaTeX Project Public
%% License, either version 1.2 of this license or (at your option) any
%% later version.  The latest version of this license is in
%%    http://www.latex-project.org/lppl.txt
%% and version 1.2 or later is part of all distributions of LaTeX
%% version 1999/12/01 or later.
%%
%% The list of all files belonging to the 'CAS Bundle' is
%% given in the file `manifest.txt'.
%%
%% Template article for cas-sc documentclass for
%% double column output.

\documentclass[a4paper,fleqn]{cas-sc}
\usepackage{njfutitlepage}
\usepackage[authoryear]{natbib}
\usepackage{csquotes}
\usepackage[UTF8]{ctex}
\usepackage{gbt7714}
\usepackage{makecell}
\usepackage{graphicx}
\usepackage{booktabs}
\usepackage{multirow}
\usepackage{url}
\usepackage{listings}
\usepackage{xcolor}
\usepackage{hyperref}
\usepackage{amsmath}
\usepackage{amssymb}
\usepackage{placeins}

\newcommand{\refSec}[1]{\hyperref[#1]{Section~\ref*{#1}}}
\newcommand{\refFig}[1]{\hyperref[#1]{Fig.~\ref*{#1}}}
\newcommand{\refApp}[1]{\hyperref[#1]{Appendix~\ref*{#1}}}
\renewcommand{\figurename}{Fig.}


\lstset{
    language=bash,
    basicstyle=\footnotesize\ttfamily,
    keywordstyle=\color{blue}\bfseries,
    commentstyle=\color{gray},
    stringstyle=\color{red},
    numbers=left,
    numberstyle=\tiny\color{gray},
    stepnumber=1,
    numbersep=8pt,
    backgroundcolor=\color{gray!10},
    frame=single,
    tabsize=2,
    captionpos=b,
    breaklines=true,
    breakatwhitespace=false,
    escapeinside={\%*}{*)}
}

%%%Author definitions
\def\tsc#1{\csdef{#1}{\textsc{\lowercase{#1}}\xspace}}
\tsc{WGM}
\tsc{QE}
\tsc{EP}
\tsc{PMS}
\tsc{BEC}
\tsc{DE}
%%%

\begin{document}
\let\WriteBookmarks\relax
\def\floatpagepagefraction{1}
\def\textpagefraction{.001}

\shortauthors{7组（王孜悠）}
\shorttitle{}
\title[mode = title]{FeHALS 航路规划与激光仿真系统 API 技术文档}

\Course{GIS设计与开发}
\Major{测绘工程}
\Grade{2023级}
\Group{7 组}

\Member[2310604116]{王 孜悠}
\Member[2310604109]{梅 康凯}
\Member[2310604115]{唐 如意}
\Member[2310604122]{杨 杰瑞}
\Member[2310604117]{温 璞}
\Member[2110604112]{刘 继明}
\Member[2110604106]{李 英杰}
\Member[2310604103]{杜 仲宇}
\Member[2310604113]{施 祺}
\Member[2310604118]{向 宇鸣}

\Teacher{隋 铭明}
\Teacher{那 嘉明}
\Teacher{陈 动}
\Teacher{朱 杰}


\MakeNJFUTitlePage[UseDocTitle]
\tableofcontents
\newpage

\section{系统概述}

\subsection{项目背景}
FeHALS（Frontend of HELIOS++ for Airborne Laser Scanning）是一个基于Web的三维可视化航路规划与激光仿真系统，旨在解决HELIOS++原生命令行交互复杂、缺乏三维可视化航路规划能力的问题。本系统为航测工程师、遥感研究人员以及相关专业师生提供一个直观易用的平台，支持从场景构建、航路设计、仿真执行到结果可视化的全流程操作。

\subsection{系统架构}
系统采用浏览器/服务器（B/S）三层架构：
\begin{itemize}
    \item \textbf{前端层}：基于Vue 3与Three.js构建，负责三维场景渲染、航点交互编辑、参数配置与结果展示
    \item \textbf{后端层}：基于FastAPI实现，负责API服务、文件生成、仿真调度与数据处理
    \item \textbf{仿真引擎层}：HELIOS++激光扫描仿真引擎，独立运行并由后端调用
\end{itemize}

\begin{figure}[pos=!htbp]
    \centering
    \caption{FeHALS系统架构与数据流}
    \includegraphics[width=0.8\textwidth]{figs/gernal_fehals_arch.pdf}
    \label{fig:architecture}
\end{figure}

\section{相关技术}\label{sec:tech}

\subsection{HELIOS++仿真框架}
HELIOS++是一款基于C++的开源激光雷达仿真框架，采用光线追踪技术对激光脉冲在三维场景中的传播过程进行建模，能够模拟机载、地面、车载与无人机等多种平台，支持全波形（full-waveform）输出与LAS、LAZ、XYZ等多种点云格式。其仿真配置以XML文件组织：场景（scene）文件定义参与扫描的三维地物模型；平台（platform）与扫描器（scanner）目录文件定义飞行平台与扫描仪的物理参数；扫描任务（survey）文件将场景、平台、扫描器与航迹（trajectory）关联起来，并设定脉冲频率、扫描频率与扫描角等核心参数。

\subsection{Web三维可视化与Three.js}
Three.js是一个基于WebGL \citep{webgl} 的开源JavaScript三维图形库，封装了场景（Scene）、相机（Camera）、渲染器（Renderer）与网格（Mesh）等核心对象，支持OBJ、GLTF、STL等多种三维模型格式的加载，并提供了轨道控制（OrbitControls）、拖拽控制（DragControls）与射线拾取（Raycaster）等交互组件。借助Three.js，可在浏览器中构建轻量级的三维可视化场景，实现模型展示、点云渲染与鼠标交互拾取，为航路规划提供直观的操作界面。

\subsection{Vue 3前端框架}
Vue 3是一个渐进式JavaScript框架，采用Composition API提供更好的代码组织和类型支持。其核心特性包括响应式系统、组件化开发和虚拟DOM。本系统前端基于Vue 3构建，使用Pinia进行状态管理，通过组件化方式组织用户界面，实现了高效的3D交互和参数配置功能。

\subsection{FastAPI与后端异步任务}
FastAPI是一个基于Python的高性能Web框架，基于类型注解与Pydantic \citep{pydantic} 实现自动的数据校验与接口文档生成，并原生支持异步编程与WebSocket。本系统后端基于FastAPI构建REST API \citep{fielding2000rest} 与WebSocket日志通道，通过asyncio子进程异步调用HELIOS++可执行文件，实时捕获其标准输出并以WebSocket推送到前端，实现仿真过程的实时监控。

\section{系统需求分析与设计}\label{sec:design}

\subsection{用户需求分析}
本系统的目标用户为航测工程师、遥感研究人员以及相关专业师生。其核心需求可归纳为三个方面：一是直观的航路规划，用户希望在不编写配置文件的情况下，通过三维场景中的鼠标交互即可完成航点设定与航线编辑；二是快速的仿真执行，用户能够通过图形化界面配置平台与扫描仪参数，一键调用HELIOS++引擎并实时查看运行状态；三是便捷的结果可视化，仿真生成的点云能够自动加载到三维场景中，并支持按高度或强度着色以及尺寸、透明度的调节。

\subsection{系统总体架构}
系统采用浏览器/服务器（B/S）架构，分为前端、后端与仿真引擎三层，其总体架构如图\ref{fig:architecture}所示。前端基于Vue 3构建，负责三维场景渲染、航点交互编辑、参数配置与结果展示，通过HTTP \citep{rfc9110,rfc9113,rfc9114} 与WebSocket与后端通信；后端基于FastAPI，负责接收前端请求，生成HELIOS++所需的航迹文件与XML配置文件，调用HELIOS++可执行文件执行仿真，并解析与返回点云结果；仿真引擎层为HELIOS++，独立于前后端，由后端以子进程方式调用。各模块间通过明确的REST API与WebSocket消息格式进行数据交换，形成完整的工作流。

\subsection{功能优先级划分}
结合开发环境与工程约束，系统功能划分为三个优先级：核心功能（MVP）包括三维场景渲染、模型加载、交互式航点编辑、航迹导出、仿真参数配置、仿真执行、日志控制台与点云渲染；重要功能包括场景要素属性编辑、自动航线生成、航高自动计算、仿真预设方案、点云特征统计与点云下载；拓展功能包括仿真进度实时反馈、点云覆盖度分析、多任务队列管理与仿真过程动画。

\section{系统设计与实现}\label{sec:impl}

\subsection{系统模块划分}
系统实现按前后端分层组织，前端基于Vue 3与Node.js \citep{nodejs} 生态构建，目录下包含组件（components）、状态管理（stores）与组合式函数（composables）三部分：组件层包含三维场景Scene3D.vue、参数配置面板ControlPanel.vue、航点列表WaypointList.vue与日志控制台LogConsole.vue；stores层基于Pinia维护场景（scene）、航点（waypoints）与仿真（simulation）三份响应式状态；composables层封装Three.js场景管理（useThreeScene）、航点交互（useWaypoints）与后端接口调用（useHeliosAPI）。

后端基于FastAPI构建，api目录下包含REST路由routes.py与WebSocket端点websocket.py，models目录定义Pydantic数据模型schemas.py，services目录封装航迹生成trajectory\_generator.py、配置生成config\_generator.py、HELIOS++调用helios\_service.py与点云解析pointcloud\_parser.py。

\subsection{三维场景渲染与航点交互}
三维场景基于Three.js构建，Scene3D.vue负责初始化场景、相机、渲染器与灯光，并添加地面网格与拾取平面。场景使用Z-up坐标系（与HELIOS++一致，z为高程轴），通过OrbitControls实现相机的旋转、平移与缩放；模型加载采用OBJLoader、GLTFLoader与STLLoader，加载完成后根据模型包围盒自动调整相机视角，并基于模型顶点包围盒推断其up轴（Y-up模型自动旋转至Z-up）。航点交互通过自绘拖拽实现：鼠标左键点击经Raycaster与地面平面求交，在交点处添加球体表示的航点（恒为z=0贴地），相邻航点之间以白色折线连接；拖拽航点时锁定相机（`controls.enabled=false`），确保不干扰视角操作；通过Delete键或右键菜单删除航点。

\subsection{前端组件与状态管理}
侧边栏采用多Tab布局，包含仿真参数、点云、模型列表、航迹与设置五个Tab。ControlPanel.vue提供仿真参数表单，分为「载体参数」与「传感器参数」两个模块，各参数旁显示有效范围，只读参数以灰色禁用显示；参数数据来源于从HELIOS++扫描器与平台目录文件中提取的规格库，切换平台类型时自动重置参数至默认值。WaypointList.vue以列表形式展示航点序号与坐标，支持单个删除与整体清空。LogConsole.vue展示分级日志（INFO/WARNING/ERROR），支持自动滚动与清空。

三个Pinia store分别维护场景模型与点云渲染参数（scene）、航点坐标数组（waypoints）以及仿真任务状态、进度、日志与结果（simulation）。前端通过useHeliosAPI.js中的Axios实例调用后端REST接口，并通过connectLogWS建立WebSocket连接以接收仿真日志。

\subsection{后端API接口设计}
后端基于FastAPI实现，提供完整的RESTful API接口，主要分为以下几类：

\begin{table}[pos=!htbp]
    \centering
    \caption{FeHALS REST API接口列表}
    \label{tab:api-interfaces}
    \begin{tabular}{@{} LLL @{}}
        \toprule
        \textbf{接口组} & \textbf{HTTP方法}                          & \textbf{功能说明} \\
        \midrule
        模型管理         & GET /api/models                          & 获取所有上传的模型列表   \\
                     & POST /api/models/upload                  & 上传3D模型文件      \\
                     & DELETE /api/models/\{id\}                & 删除指定模型        \\
        航迹管理         & POST /api/trajectory/generate            & 从航点生成航迹文件     \\
                     & GET /api/trajectories                    & 获取所有航迹列表      \\
                     & GET /api/trajectories/download/\{id\}    & 下载航迹文件        \\
        配置管理         & POST /api/config/generate                & 生成仿真配置文件      \\
        仿真执行         & POST /api/simulation/run                 & 执行HELIOS++仿真  \\
                     & GET /api/simulation/status/\{task\_id\}  & 获取仿真状态        \\
                     & GET /api/simulation/logs/\{task\_id\}    & 获取仿真日志        \\
                     & POST /api/simulation/cancel/\{task\_id\} & 取消仿真任务        \\
        结果管理         & GET /api/results/\{task\_id\}            & 获取点云数据        \\
                     & GET /api/results/\{task\_id\}/download   & 下载原始结果        \\
        环境诊断         & GET /api/env/diagnose                    & 检查系统环境        \\
        \bottomrule
    \end{tabular}
\end{table}

\subsection{航迹与配置文件生成}
trajectory\_generator.py将航点数组按添加顺序、以1秒时间间隔生成HELIOS++原生航迹文件（.trj，CSV格式，列顺序为时间、roll、pitch、yaw、x、y、z），所有航点统一为恒定飞行高度（仅用x、y定义水平路径）。系统还支持弓字形自动航迹生成：用户在三维场景中点击两个角点定义矩形区域，系统根据设定航线间距自动生成往返覆盖的全覆盖航线。

config\_generator.py负责生成HELIOS++所需的两个XML文件：场景文件通过objloader滤镜引用三维模型，并指定模型的up轴（Y-up自动旋转至Z-up）；扫描任务（survey）文件引用场景、平台与扫描器目录，并将平台类型映射为对应的平台目录项（UAV与Airborne均使用`copter\_linearpath`平台、`riegl\_vux-1uav`扫描器），将脉冲频率（kHz）与扫描角（±半角）换算为HELIOS++所需的脉冲频率（Hz）与总扫描角。

\subsection{HELIOS++调用与点云渲染}
helios\_service.py维护仿真任务注册表，通过asyncio.create\_subprocess\_exec以子进程方式调用helios++可执行文件，指定资源目录（--assets）、输出目录（--output）与输出格式参数，实时捕获其标准输出/错误流并解析进度百分比，经WebSocket推送至前端，同时支持超时终止与任务取消。

仿真结束后，pointcloud\_parser.py利用laspy \citep{laspy} 解析LAS/LAZ点云或按文本方式解析XYZ点云，提取坐标与强度、计算空间范围，并对超过百万级的点云进行降采样。前端将返回的点云坐标构建为BufferGeometry与Points对象，支持按高度（蓝色至红色渐变）、按强度或固定颜色着色，并可实时调节点尺寸与透明度。

\subsection{项目组织结构}
\begin{figure}[pos=!htbp]
    \centering
    \caption{FeHALS项目目录结构}
    \begin{lstlisting}[language=bash]
FeHALS/
├── backend/                 # 后端代码
│   ├── app/
│   │   ├── main.py         # FastAPI应用入口
│   │   ├── config.py       # 配置管理
│   │   ├── api/            # API路由
│   │   ├── models/         # 数据模型
│   │   ├── services/       # 业务服务
│   │   └── static/         # 静态文件目录
│   ├── requirements.txt    # Python依赖
│   └── run.py             # 后端启动脚本
├── frontend/              # 前端代码
│   ├── src/
│   │   ├── components/    # Vue组件
│   │   ├── composables/  # 组合式函数
│   │   ├── stores/       # Pinia状态管理
│   │   └── assets/       # 静态资源
│   ├── package.json
│   └── vite.config.js
├── doc/                   # 文档目录
│   ├── src/
│   │   ├── DESCRIPTION.tex
│   │   ├── DESIGNBOOK.tex
│   │   └── DOCUMENT.tex
│   └── Makefile           # 构建脚本
└── run.sh                 # Unix启动脚本
  \end{lstlisting}
    \label{fig:project-structure}
\end{figure}

\section{部署与运行指南}

\subsection{环境要求}
\begin{table}[pos=!htbp]
    \centering
    \caption{系统环境要求}
    \begin{tabular}{@{} ll @{}}
        \toprule
        \textbf{组件} & \textbf{要求}      \\
        \midrule
        Node.js     & >= 16.0.0        \\
        Python      & >= 3.8.0         \\
        npm         & >= 7.0.0         \\
        LaTeX       & XeLaTeX (用于文档编译) \\
        \bottomrule
    \end{tabular}
\end{table}

\subsection{安装步骤}
\begin{enumerate}
    \item \textbf{后端安装}
          \begin{lstlisting}[language=bash]
  # 1. 进入后端目录
  cd backend

  # 2. 创建虚拟环境
  python -m venv venv

  # 3. 激活虚拟环境
  # Linux/Mac
  source venv/bin/activate
  # Windows
  venv\Scripts\activate

  # 4. 安装依赖
  pip install -r requirements.txt

  # 5. 设置环境变量
  export HELIOS_PATH=/path/to/helios++
  export HELIOS_REPO=/path/to/helios++/repo
  export HELIOS_ASSETS=/path/to/assets
  \end{lstlisting}

    \item \textbf{前端安装}
          \begin{lstlisting}[language=bash]
  # 1. 进入前端目录
  cd frontend

  # 2. 安装依赖
  npm install

  # 3. 启动开发服务器
  npm run dev
  \end{lstlisting}
\end{enumerate}

\subsection{启动脚本}
项目提供了启动脚本用于快速启动服务：
\begin{lstlisting}[language=bash,caption=run.sh脚本]
#!/bin/bash

# 停止所有服务
stop_services() {
    echo "停止所有服务..."
    pkill -f "uvicorn" || true
    pkill -f "vite" || true
}

# 启动后端
start_backend() {
    echo "启动后端服务..."
    cd backend
    source venv/bin/activate
    python -m uvicorn app.main:app --reload --host 0.0.0.0 --port 8000 &
    cd ..
}

# 启动前端
start_frontend() {
    echo "启动前端服务..."
    cd frontend
    npm run dev &
    cd ..
}

# 主程序
case "$1" in
    start)
        start_backend
        start_frontend
        echo "服务已启动"
        echo "前端: http://localhost:3000"
        echo "后端: http://localhost:8000"
        echo "API文档: http://localhost:8000/docs"
        ;;
    stop)
        stop_services
        echo "服务已停止"
        ;;
    *)
        echo "用法: $0 {start|stop}"
        exit 1
        ;;
esac
\end{lstlisting}

\subsection{配置说明}
\begin{itemize}
    \item \textbf{后端配置}
          \begin{lstlisting}[language=bash]
  # 环境变量示例
  export FEHALS_HOST=0.0.0.0
  export FEHALS_PORT=8000
  export HELIOS_PATH=/usr/local/bin/helios++
  export HELIOS_REPO=/opt/helios++
  export HELIOS_ASSETS=/opt/helios++/assets
  export FEHALS_SIM_TIMEOUT=300
  export FEHALS_CORS_ORIGINS=http://localhost:3000,http://127.0.0.1:3000
  \end{lstlisting}

    \item \textbf{前端配置}
          \begin{lstlisting}[language=bash]
  // vite.config.js
  export default defineConfig({
    server: {
      proxy: {
        '/api': {
          target: 'http://localhost:8000',
          changeOrigin: true
        },
        '/ws': {
          target: 'ws://localhost:8000',
          ws: true
        }
      }
    }
  })
  \end{lstlisting}
\end{itemize}

\section{API接口详解}\label{sec:api}

\subsection{REST API接口}

\subsubsection{模型管理接口}
\begin{table}[pos=!htbp]
    \centering
    \caption{模型管理API接口}
    \label{tab:api-models}
    \begin{tabular}{@{} LLL @{}}
        \toprule
        \textbf{接口路径}                    & \textbf{HTTP方法} & \textbf{功能说明}              \\
        \midrule
        \texttt{/api/models}             & GET             & 获取所有已上传的模型列表               \\
        \texttt{/api/models/upload}      & POST            & 上传3D模型文件（支持OBJ、GLTF、STL格式） \\
        \texttt{/api/models/\{id\}}      & DELETE          & 删除指定ID的模型                  \\
        \texttt{/api/models/\{id\}/info} & GET             & 获取指定模型的详细信息                \\
        \texttt{/api/models/\{id\}/bbox} & GET             & 获取指定模型的包围盒信息               \\
        \bottomrule
    \end{tabular}
\end{table}

\paragraph{上传模型接口示例}
\begin{lstlisting}[language=bash,caption=上传3D模型请求]
# 请求示例（使用curl）
curl -X POST \
  http://localhost:8000/api/models/upload \
  -H "Content-Type: multipart/form-data" \
  -F "file=@/path/to/model.obj" \
  -F "name=test_model" \
  -F "description=测试模型文件"

# 响应示例（JSON）
{
  "id": "model_123",
  "name": "test_model",
  "description": "测试模型文件",
  "filename": "model.obj",
  "size": 1024000,
  "format": "obj",
  "upload_time": "2024-01-01T12:00:00Z",
  "bbox": {
    "min": {"x": -10, "y": -10, "z": 0},
    "max": {"x": 10, "y": 10, "z": 20}
  }
}
\end{lstlisting}

\subsubsection{航迹管理接口}
\begin{table}[pos=!htbp]
    \centering
    \caption{航迹管理API接口}
    \label{tab:api-trajectory}
    \begin{tabular}{@{} LLL @{}}
        \toprule
        \textbf{接口路径}                              & \textbf{HTTP方法} & \textbf{功能说明}       \\
        \midrule
        \texttt{/api/trajectory/generate}          & POST            & 从航点列表生成HELIOS++航迹文件 \\
        \texttt{/api/trajectories}                 & GET             & 获取所有保存的航迹列表         \\
        \texttt{/api/trajectories/\{id\}}          & GET             & 获取指定航迹的详细信息         \\
        \texttt{/api/trajectories/\{id\}/download} & GET             & 下载指定航迹的.trj文件       \\
        \texttt{/api/trajectories/\{id\}/preview}  & GET             & 获取航迹的3D预览数据         \\
        \bottomrule
    \end{tabular}
\end{table}

\paragraph{生成航迹接口示例}
\begin{lstlisting}[language=bash,caption=生成航迹请求]
# 请求体（JSON）
{
  "waypoints": [
    {"id": 1, "x": 0, "y": 0, "z": 100},
    {"id": 2, "x": 100, "y": 0, "z": 100},
    {"id": 3, "x": 100, "y": 100, "z": 100},
    {"id": 4, "x": 0, "y": 100, "z": 100}
  ],
  "flight_height": 100,
  "time_interval": 1.0,
  "pattern": "manual"
}

# 生成的航迹文件格式（CSV示例）
time,roll,pitch,yaw,x,y,z
0.0,0.0,0.0,0.0,0.0,0.0,100.0
1.0,0.0,0.0,0.0,100.0,0.0,100.0
2.0,0.0,0.0,0.0,100.0,100.0,100.0
3.0,0.0,0.0,0.0,0.0,100.0,100.0
\end{lstlisting}

\subsubsection{仿真执行接口}
\begin{table}[pos=!htbp]
    \centering
    \caption{仿真执行API接口}
    \label{tab:api-simulation}
    \begin{tabular}{@{} LLL @{}}
        \toprule
        \textbf{接口路径}                                 & \textbf{HTTP方法} & \textbf{功能说明}  \\
        \midrule
        \texttt{/api/simulation/run}                  & POST            & 启动HELIOS++仿真任务 \\
        \texttt{/api/simulation/status/\{task\_id\}}  & GET             & 获取仿真任务状态       \\
        \texttt{/api/simulation/logs/\{task\_id\}}    & GET             & 获取仿真任务日志       \\
        \texttt{/api/simulation/cancel/\{task\_id\}}  & POST            & 取消正在执行的仿真任务    \\
        \texttt{/api/simulation/restart/\{task\_id\}} & POST            & 重启失败的仿真任务      \\
        \texttt{/api/simulation/history}              & GET             & 获取历史仿真任务列表     \\
        \bottomrule
    \end{tabular}
\end{table}

\paragraph{启动仿真接口示例}
\begin{lstlisting}[language=bash,caption=启动仿真请求]
# 请求体（JSON）
{
  "survey_config": {
    "platform": "copter_linearpath",
    "scanner": "riegl_vux-1uav",
    "trajectory_id": "trajectory_123",
    "scene_id": "scene_456"
  },
  "simulation_params": {
    "pulse_frequency": 100000,
    "scan_frequency": 100,
    "scan_angle": 15,
    "flight_speed": 30,
    "output_format": "LAS"
  },
  "output_settings": {
    "output_path": "/tmp/simulation_results",
    "file_prefix": "simulation_001",
    "include_intensity": true,
    "include_fullwave": false
  }
}

# 任务状态响应
{
  "task_id": "sim_789",
  "status": "running",
  "progress": 0.45,
  "start_time": "2024-01-01T12:00:00Z",
  "estimated_completion": "2024-01-01T12:05:00Z"
}
\end{lstlisting}

\subsubsection{结果管理接口}
\begin{table}[pos=!htbp]
    \centering
    \caption{结果管理API接口}
    \label{tab:api-results}
    \begin{tabular}{@{} LLL @{}}
        \toprule
        \textbf{接口路径}                                 & \textbf{HTTP方法} & \textbf{功能说明}    \\
        \midrule
        \texttt{/api/results/\{task\_id\}}            & GET             & 获取仿真结果概览信息       \\
        \texttt{/api/results/\{task\_id\}/download}   & GET             & 下载完整的仿真结果文件      \\
        \texttt{/api/results/\{task\_id\}/stats}      & GET             & 获取点云统计数据         \\
        \texttt{/api/results/\{task\_id\}/preview}    & GET             & 获取点云预览数据（前1000点） \\
        \texttt{/api/results/\{task\_id\}/downsample} & POST            & 生成降采样后的点云数据      \\
        \texttt{/api/results/\{task\_id\}/export}     & POST            & 导出为指定格式的点云文件     \\
        \bottomrule
    \end{tabular}
\end{table}

\subsection{WebSocket接口}

\subsubsection{日志推送}
\begin{table}[pos=!htbp]
    \centering
    \caption{WebSocket日志消息格式}
    \label{tab:ws-logs}
    \begin{tabular}{@{} LLL @{}}
        \toprule
        \textbf{消息类型}   & \textbf{数据字段}         & \textbf{说明}                              \\
        \midrule
        \texttt{log}    & \texttt{timestamp}    & 日志时间戳                                    \\
                        & \texttt{level}        & 日志级别（INFO/WARNING/ERROR）                 \\
                        & \texttt{message}      & 日志消息内容                                   \\
                        & \texttt{progress}     & 进度百分比（0-100）                             \\
        \midrule
        \texttt{status} & \texttt{task\_id}     & 任务ID                                     \\
                        & \texttt{status}       & 任务状态（running/completed/failed/cancelled） \\
                        & \texttt{progress}     & 当前进度                                     \\
                        & \texttt{error}        & 错误信息（如果有）                                \\
        \midrule
        \texttt{result} & \texttt{task\_id}     & 任务ID                                     \\
                        & \texttt{result\_path} & 结果文件路径                                   \\
                        & \texttt{point\_count} & 点云总数                                     \\
                        & \texttt{file\_size}   & 结果文件大小                                   \\
        \bottomrule
    \end{tabular}
\end{table}

\paragraph{WebSocket连接示例}
\begin{lstlisting}[caption=WebSocket客户端连接]
// 建立WebSocket连接
const ws = new WebSocket('ws://localhost:8000/ws/simulation/sim_789');

// 接收日志消息
ws.onmessage = function(event) {
    const data = JSON.parse(event.data);

    switch(data.type) {
        case 'log':
            console.log(`[${data.level}] ${data.timestamp}: ${data.message}`);
            updateProgress(data.progress);
            break;
        case 'status':
            console.log(`Task ${data.task_id} status: ${data.status}`);
            break;
        case 'result':
            console.log(`Simulation completed: ${data.result_path}`);
            downloadResults(data.task_id);
            break;
    }
};

// 发送控制消息
function cancelSimulation() {
    ws.send(JSON.stringify({
        type: 'command',
        action: 'cancel'
    }));
}
\end{lstlisting}

\subsection{数据模型定义}

\subsubsection{航点数据模型}
\begin{lstlisting}[language=python,caption=Waypoint Pydantic模型]
# models/schemas.py
from pydantic import BaseModel
from typing import List, Optional
from datetime import datetime

class Waypoint(BaseModel):
    id: int
    x: float
    y: float
    z: float = 100.0
    timestamp: Optional[datetime] = None

class WaypointList(BaseModel):
    waypoints: List[Waypoint]
    name: Optional[str] = "default_trajectory"
    created_at: datetime = datetime.now()
\end{lstlisting}

\subsubsection{仿真参数数据模型}
\begin{lstlisting}[language=python,caption=SimulationParams Pydantic模型]
class PlatformConfig(BaseModel):
    type: str  # "copter_linearpath", "fixedwing", "ground"
    name: str
    speed: float
    min_height: float
    max_height: float

class ScannerConfig(BaseModel):
    type: str  # "riegl_vux-1uav", "livox-mid70", "custom"
    pulse_frequency: int  # Hz
    scan_frequency: float  # Hz
    scan_angle: float  # degrees
    scan_direction: str  # "nadir", "oblique"

class SimulationConfig(BaseModel):
    platform: PlatformConfig
    scanner: ScannerConfig
    trajectory_id: str
    output_format: str = "LAS"
    output_path: str
    timeout: int = 300  # seconds

# 全局配置常量（实际应从config.py导入）
HELLOS_PATH = "/usr/local/bin/helios++"
HELLOS_REPO = "/opt/helios++/repo"
HELLOS_ASSETS = "/opt/helios++/assets"
MAX_FILE_SIZE = 100 * 1024 * 1024  # 100MB
ALLOWED_EXTENSIONS = [".obj", ".gltf", ".glb", ".stl"]
\end{lstlisting}

\subsubsection{任务状态数据模型}
\begin{lstlisting}[language=python,caption=TaskStatus Pydantic模型]
class TaskStatus(str, Enum):
    PENDING = "pending"
    RUNNING = "running"
    COMPLETED = "completed"
    FAILED = "failed"
    CANCELLED = "cancelled"

class SimulationTask(BaseModel):
    task_id: str
    status: TaskStatus
    progress: float = 0.0
    created_at: datetime
    started_at: Optional[datetime] = None
    completed_at: Optional[datetime] = None
    error_message: Optional[str] = None
    config: SimulationConfig
\end{lstlisting}

\section{前端实现详解}\label{sec:frontend}

\subsection{Vue 3应用架构}

\subsubsection{组件结构}
前端采用Vue 3 Composition API构建，主要组件包括：
\begin{lstlisting}[language=bash,caption=前端组件结构]
frontend/src/
├── components/          # Vue组件
│   ├── Scene3D.vue     # 3D场景渲染组件
│   ├── ControlPanel.vue  # 参数配置面板
│   ├── WaypointList.vue  # 航点列表组件
│   ├── LogConsole.vue    # 日志控制台
│   └── ModelManager.vue  # 模型管理组件
├── stores/            # Pinia状态管理
│   ├── scene.ts       # 场景状态
│   ├── waypoints.ts   # 航点状态
│   └── simulation.ts  # 仿真状态
└── composables/      # 组合式函数
    ├── useThreeScene.ts  # Three.js场景管理
    ├── useWaypoints.ts   # 航点交互逻辑
    └── useHeliosAPI.ts   # API调用封装
\end{lstlisting}

\subsubsection{3D场景组件实现}
\begin{lstlisting}[caption=Scene3D.vue组件核心代码]
// components/Scene3D.vue
import { ref, onMounted, onBeforeUnmount } from 'vue'
import * as THREE from 'three'
import { OrbitControls } from 'three/examples/jsm/controls/OrbitControls'
import { OBJLoader, GLTFLoader, STLLoader } from 'three/examples/jsm/loaders'

export default {
  setup() {
    const container = ref(null)
    const scene = ref(null)
    const camera = ref(null)
    const renderer = ref(null)
    const controls = ref(null)

    // 初始化3D场景
    const initScene = () => {
      // 创建场景
      scene.value = new THREE.Scene()
      scene.value.background = new THREE.Color(0xf0f0f0)

      // 设置相机
      camera.value = new THREE.PerspectiveCamera(
        75,
        container.value.clientWidth / container.value.clientHeight,
        0.1,
        1000
      )
      camera.value.position.set(50, 50, 50)

      // 创建渲染器
      renderer.value = new THREE.WebGLRenderer({ antialias: true })
      renderer.value.setSize(
        container.value.clientWidth,
        container.value.clientHeight
      )
      container.value.appendChild(renderer.value.domElement)

      // 添加轨道控制
      controls.value = new OrbitControls(camera.value, renderer.value.domElement)
      controls.value.enableDamping = true

      // 添加灯光
      const ambientLight = new THREE.AmbientLight(0xffffff, 0.6)
      scene.value.add(ambientLight)

      const directionalLight = new THREE.DirectionalLight(0xffffff, 0.4)
      directionalLight.position.set(100, 100, 100)
      scene.value.add(directionalLight)

      // 添加地面
      const groundGeometry = new THREE.PlaneGeometry(1000, 1000)
      const groundMaterial = new THREE.MeshLambertMaterial({
        color: 0xcccccc,
        side: THREE.DoubleSide
      })
      const ground = new THREE.Mesh(groundGeometry, groundMaterial)
      ground.rotation.x = -Math.PI / 2
      scene.value.add(ground)

      // 开始渲染循环
      animate()
    }

    // 动画循环
    const animate = () => {
      requestAnimationFrame(animate)
      controls.value.update()
      renderer.value.render(scene.value, camera.value)
    }

    // 处理窗口大小变化
    const handleResize = () => {
      camera.value.aspect = container.value.clientWidth / container.value.clientHeight
      camera.value.updateProjectionMatrix()
      renderer.value.setSize(
        container.value.clientWidth,
        container.value.clientHeight
      )
    }

    onMounted(() => {
      initScene()
      window.addEventListener('resize', handleResize)
    })

    onBeforeUnmount(() => {
      window.removeEventListener('resize', handleResize)
    })

    return {
      container,
      addWaypoint,
      clearWaypoints
    }
  }
}
\end{lstlisting}

\subsection{Three.js集成与3D交互}

\subsubsection{模型加载系统}
\begin{lstlisting}[caption=模型加载组合式函数]
// composables/useModelLoader.ts
import { ref } from 'vue'
import * as THREE from 'three'
import { OBJLoader } from 'three/examples/jsm/loaders/OBJLoader'
import { GLTFLoader } from 'three/examples/jsm/loaders/GLTFLoader'
import { STLLoader } from 'three/examples/jsm/loaders/STLLoader'

export function useModelLoader() {
  const models = ref<Map<string, THREE.Object3D>>(new Map())
  const loading = ref(false)
  const error = ref<string | null>(null)

  const loadModel = async (url: string, name: string, format: string) => {
    loading.value = true
    error.value = null

    try {
      let loader: any
      let object: THREE.Object3D

      switch (format.toLowerCase()) {
        case 'obj':
          loader = new OBJLoader()
          object = await loader.loadAsync(url)
          break
        case 'gltf':
        case 'glb':
          loader = new GLTFLoader()
          const gltf = await loader.loadAsync(url)
          object = gltf.scene
          break
        case 'stl':
          loader = new STLLoader()
          const geometry = await loader.loadAsync(url)
          const material = new THREE.MeshStandardMaterial({
            color: 0x88ccff,
            wireframe: false
          })
          object = new THREE.Mesh(geometry, material)
          break
        default:
          throw new Error(`Unsupported format: ${format}`)
      }

      // 自动调整模型朝向（Y-up转Z-up）
      const bbox = new THREE.Box3().setFromObject(object)
      const center = bbox.getCenter(new THREE.Vector3())
      const size = bbox.getSize(new THREE.Vector3())

      // 计算模型的主朝向
      const maxAxis = size.x > size.y ? 'x' : 'y'
      const rotationAngle = maxAxis === 'x' ? Math.PI / 2 : 0

      // 旋转模型
      object.rotation.z = rotationAngle
      object.position.copy(center.multiplyScalar(-1))

      // 计算包围盒
      const worldBbox = new THREE.Box3().setFromObject(object)
      const worldCenter = worldBbox.getCenter(new THREE.Vector3())
      const worldSize = worldBbox.getSize(new THREE.Vector3())

      models.value.set(name, object)

      return {
        object,
        bbox: {
          min: worldBbox.min,
          max: worldBbox.max,
          center: worldCenter,
          size: worldSize
        }
      }
    } catch (err) {
      error.value = err instanceof Error ? err.message : 'Failed to load model'
      throw err
    } finally {
      loading.value = false
    }
  }

  return {
    models,
    loading,
    error,
    loadModel
  }
}
\end{lstlisting}

\subsubsection{航点交互系统}
\begin{lstlisting}[caption=航点交互组合式函数]
// composables/useWaypoints.ts
import { ref, computed } from 'vue'
import * as THREE from 'three'

interface Waypoint {
  id: number
  position: THREE.Vector3
  mesh: THREE.Mesh
}

export function useWaypoints() {
  const waypoints = ref<Waypoint[]>([])
  const selectedWaypoint = ref<Waypoint | null>(null)
  const isDragging = ref(false)

  // 添加航点
  const addWaypoint = (position: THREE.Vector3) => {
    const id = Date.now()
    const geometry = new THREE.SphereGeometry(2, 32, 16)
    const material = new THREE.MeshLambertMaterial({
      color: 0xff4444,
      transparent: true,
      opacity: 0.8
    })

    const mesh = new THREE.Mesh(geometry, material)
    mesh.position.copy(position)
    mesh.position.z = 0 // 贴地
    mesh.userData = { id }

    const waypoint: Waypoint = {
      id,
      position,
      mesh
    }

    waypoints.value.push(waypoint)
    updateConnectionLines()

    return waypoint
  }

  // 删除航点
  const removeWaypoint = (id: number) => {
    const index = waypoints.value.findIndex(wp => wp.id === id)
    if (index !== -1) {
      const waypoint = waypoints.value[index]
      waypoint.mesh.removeFromParent()
      waypoints.value.splice(index, 1)
      updateConnectionLines()
    }
  }

  // 更新航点位置
  const updateWaypointPosition = (id: number, position: THREE.Vector3) => {
    const waypoint = waypoints.value.find(wp => wp.id === id)
    if (waypoint) {
      waypoint.position.copy(position)
      waypoint.mesh.position.copy(position)
      waypoint.mesh.position.z = 0
      updateConnectionLines()
    }
  }

  // 更新连接线
  const updateConnectionLines = () => {
    // 移除旧的连接线
    const oldLines = scene.value.getObjectByName('waypoint-lines')
    if (oldLines) {
      scene.value.remove(oldLines)
    }

    // 创建新的连接线
    if (waypoints.value.length > 1) {
      const points = waypoints.value.map(wp => wp.position.clone())
      const geometry = new THREE.BufferGeometry().setFromPoints(points)
      const material = new THREE.LineBasicMaterial({
        color: 0xffffff,
        linewidth: 2
      })

      const line = new THREE.Line(geometry, material)
      line.name = 'waypoint-lines'
      scene.value.add(line)
    }
  }

  // 拖拽处理
  const handleMouseDown = (event: MouseEvent) => {
    // 鼠标射线检测
    const raycaster = new THREE.Raycaster()
    const mouse = new THREE.Vector2()

    mouse.x = (event.clientX / window.innerWidth) * 2 - 1
    mouse.y = -(event.clientY / window.innerHeight) * 2 + 1

    raycaster.setFromCamera(mouse, camera.value)

    const intersects = raycaster.intersectObjects(
      waypoints.value.map(wp => wp.mesh)
    )

    if (intersects.length > 0) {
      selectedWaypoint.value = waypoints.value.find(
        wp => wp.mesh === intersects[0].object
      ) || null
      isDragging.value = true
      controls.value.enabled = false
    }
  }

  const handleMouseMove = (event: MouseEvent) => {
    if (isDragging.value && selectedWaypoint.value) {
      // 计算与地面的交点
      const raycaster = new THREE.Raycaster()
      const mouse = new THREE.Vector2()

      mouse.x = (event.clientX / window.innerWidth) * 2 - 1
      mouse.y = -(event.clientY / window.innerHeight) * 2 + 1

      raycaster.setFromCamera(mouse, camera.value)

      // 创建地面平面
      const plane = new THREE.Plane(new THREE.Vector3(0, 0, 1), 0)
      const intersectPoint = new THREE.Vector3()

      raycaster.ray.intersectPlane(plane, intersectPoint)

      // 更新航点位置
      updateWaypointPosition(selectedWaypoint.value.id, intersectPoint)
    }
  }

  const handleMouseUp = () => {
    isDragging.value = false
    selectedWaypoint.value = null
    controls.value.enabled = true
  }

  return {
    waypoints: computed(() => waypoints.value),
    addWaypoint,
    removeWaypoint,
    updateWaypointPosition,
    handleMouseDown,
    handleMouseMove,
    handleMouseUp
  }
}
\end{lstlisting}

\section{后端实现详解}\label{sec:backend}

\subsection{FastAPI应用架构}

\subsubsection{主应用配置}
\begin{lstlisting}[language=python,caption=main.py - FastAPI应用入口]
# app/main.py
from fastapi import FastAPI, HTTPException
from fastapi.middleware.cors import CORSMiddleware
from fastapi.staticfiles import StaticFiles
import uvicorn
import os
from .config import settings
from .api import router as api_router
from .api.websocket import websocket_router

app = FastAPI(
    title="FeHALS API",
    description="FeHALS航路规划与激光仿真系统API",
    version="1.0.0",
    docs_url="/docs",
    redoc_url="/redoc"
)

# 配置CORS
app.add_middleware(
    CORSMiddleware,
    allow_origins=settings.CORS_ORIGINS,
    allow_credentials=True,
    allow_methods=["*"],
    allow_headers=["*"],
)

# 挂载静态文件
app.mount("/static", StaticFiles(directory="static"), name="static")

# 注册路由
app.include_router(api_router, prefix="/api")
app.include_router(websocket_router, prefix="/ws")

@app.on_event("startup")
async def startup_event():
    """应用启动时的初始化操作"""
    # 检查HELIOS++环境
    from ..services.helios_service import check_helios_environment
    await check_helios_environment()

@app.get("/")
async def root():
    """API根路径"""
    return {
        "message": "FeHALS API",
        "version": "1.0.0",
        "docs_url": "/docs",
        "status": "running"
    }

if __name__ == "__main__":
    uvicorn.run(
        "app.main:app",
        host=settings.HOST,
        port=settings.PORT,
        reload=True
    )
\end{lstlisting}

\subsubsection{配置管理系统}
\begin{lstlisting}[language=python,caption=config.py - 应用配置]
# app/config.py
from pydantic_settings import BaseSettings
from typing import List
import os

class Settings(BaseSettings):
    # 服务器配置
    HOST: str = "0.0.0.0"
    PORT: int = 8000

    # CORS配置
    CORS_ORIGINS: List[str] = ["http://localhost:3000", "http://127.0.0.1:3000"]

    # HELIOS++路径配置
    HELIOS_PATH: str = "/usr/local/bin/helios++"
    HELIOS_REPO: str = "/opt/helios++/repo"
    HELIOS_ASSETS: str = "/opt/helios++/assets"

    # 仿真配置
    SIMULATION_TIMEOUT: int = 300  # 秒
    MAX_POINTCLOUD_SIZE: int = 10000000  # 点云最大点数
    DEFAULT_OUTPUT_FORMAT: str = "LAS"

    # 文件上传配置
    UPLOAD_DIR: str = "uploads"
    MAX_FILE_SIZE: int = 100 * 1024 * 1024  # 100MB
    ALLOWED_EXTENSIONS: List[str] = [".obj", ".gltf", ".glb", ".stl"]

    # 缓存配置
    CACHE_DIR: str = "cache"
    CACHE_TTL: int = 3600  # 1小时

    class Config:
        env_file = ".env"
        case_sensitive = True

# 全局配置实例
settings = Settings()
\end{lstlisting}

\subsection{API路由实现}

\subsubsection{REST路由}
\begin{lstlisting}[language=python,caption=api/routes.py - REST路由定义]
# api/routes.py
from fastapi import APIRouter, HTTPException, UploadFile, File, Depends
from typing import List
from ...services import (
    model_service,
    trajectory_service,
    simulation_service,
    result_service
)
from ...models.schemas import *

router = APIRouter()

@router.get("/models")
async def get_models():
    """获取所有模型列表"""
    return await model_service.get_all_models()

@router.post("/models/upload")
async def upload_model(
    file: UploadFile = File(...),
    name: str = None,
    description: str = None
):
    """上传3D模型文件"""
    # 验证文件类型
    if not file.filename.lower().endswith(tuple(ALLOWED_EXTENSIONS)):
        raise HTTPException(
            status_code=400,
            detail=f"文件类型不支持。支持的格式: {', '.join(ALLOWED_EXTENSIONS)}"
        )

    # 验证文件大小
    if file.size > MAX_FILE_SIZE:
        raise HTTPException(
            status_code=400,
            detail=f"文件大小超过限制（{MAX_FILE_SIZE / 1024 / 1024}MB）"
        )

    try:
        model = await model_service.upload_model(
            file=file,
            name=name or file.filename,
            description=description
        )
        return {"success": True, "data": model}
    except Exception as e:
        raise HTTPException(status_code=500, detail=str(e))

@router.delete("/models/{model_id}")
async def delete_model(model_id: str):
    """删除指定模型"""
    try:
        await model_service.delete_model(model_id)
        return {"success": True}
    except FileNotFoundError:
        raise HTTPException(status_code=404, detail="模型不存在")
    except Exception as e:
        raise HTTPException(status_code=500, detail=str(e))

@router.post("/trajectory/generate")
async def generate_trajectory(waypoints: WaypointList):
    """从航点生成航迹文件"""
    try:
        trajectory = await trajectory_service.generate_trajectory(waypoints)
        return {
            "success": True,
            "data": trajectory,
            "download_url": f"/api/trajectories/{trajectory.id}/download"
        }
    except Exception as e:
        raise HTTPException(status_code=500, detail=str(e))

@router.post("/simulation/run")
async def run_simulation(config: SimulationConfig):
    """启动仿真任务"""
    try:
        task = await simulation_service.run_simulation(config)
        return {
            "success": True,
            "data": task,
            "ws_url": f"ws://{settings.HOST}:{settings.PORT}/ws/simulation/{task.task_id}"
        }
    except Exception as e:
        raise HTTPException(status_code=500, detail=str(e))
\end{lstlisting}

\subsection{业务服务层}

\subsubsection{航迹生成服务}
\begin{lstlisting}[language=python,caption=services/trajectory\_generator.py - 航迹生成]
# services/trajectory_generator.py
import csv
import os
from typing import List, Dict, Any
from ..models.schemas import Waypoint, WaypointList, Trajectory

class TrajectoryGenerator:
    @staticmethod
    def generate_trajectory(waypoints: WaypointList,
                          flight_height: float = 100.0,
                          time_interval: float = 1.0) -> Trajectory:
        """从航点列表生成HELIOS++航迹文件"""

        # 生成航迹数据点
        trajectory_data = []

        for i, waypoint in enumerate(waypoints.waypoints):
            # 计算时间戳（每秒间隔）
            timestamp = i * time_interval

            # 添加航迹点
            trajectory_data.append({
                'time': timestamp,
                'roll': 0.0,
                'pitch': 0.0,
                'yaw': 0.0,
                'x': waypoint.x,
                'y': waypoint.y,
                'z': flight_height
            })

        # 生成航迹文件内容（CSV格式）
        trajectory_content = "time,roll,pitch,yaw,x,y,z\n"
        for point in trajectory_data:
            trajectory_content += (
                f"{point['time']},{point['roll']},{point['pitch']},"
                f"{point['yaw']},{point['x']},{point['y']},{point['z']}\n"
            )

        # 创建轨迹对象
        trajectory = Trajectory(
            id=f"traj_{os.urandom(4).hex()}",
            name=waypoints.name,
            waypoints=waypoints.waypoints,
            flight_height=flight_height,
            time_interval=time_interval,
            data=trajectory_data,
            content=trajectory_content
        )

        return trajectory

    @staticmethod
    def generate_grid_trajectory(
        bounds: Dict[str, float],
        line_spacing: float,
        flight_height: float = 100.0,
        time_interval: float = 1.0
    ) -> Trajectory:
        """生成弓字形（网格）航迹"""

        # 计算网格点
        x_min, y_min = bounds['min_x'], bounds['min_y']
        x_max, y_max = bounds['max_x'], bounds['max_y']

        points = []

        # 生成往返航点
        y_line = y_min
        direction = 1  # 1: 从左到右, -1: 从右到左

        while y_line <= y_max:
            if direction == 1:
                x_range = range(int(x_min), int(x_max), int(line_spacing))
            else:
                x_range = range(int(x_max), int(x_min), -int(line_spacing))

            for x in x_range:
                points.append({
                    'x': float(x),
                    'y': float(y_line),
                    'z': flight_height
                })

            # 下一行
            y_line += line_spacing
            direction *= -1

        # 创建航点列表
        waypoint_list = WaypointList(
            waypoints=[
                Waypoint(
                    id=i+1,
                    x=p['x'],
                    y=p['y'],
                    z=p['z']
                ) for i, p in enumerate(points)
            ],
            name="grid_trajectory"
        )

        # 生成航迹
        return TrajectoryGenerator.generate_trajectory(
            waypoint_list,
            flight_height,
            time_interval
        )
\end{lstlisting}

\subsubsection{点云处理服务}
\begin{lstlisting}[language=python,caption=services/pointcloud\_parser.py - 点云处理]
# services/pointcloud_parser.py
import laspy
import numpy as np
import os
from typing import Dict, Any, List, Optional
from ..models.schemas import PointCloudStats, PointCloudData

class PointCloudParser:
    @staticmethod
    def parse_las_laz(file_path: str) -> Dict[str, Any]:
        """解析LAS/LAZ格式点云"""

        try:
            with laspy.open(file_path) as las_file:
                las = las_file.read()

                # 提取点云数据
                points = {
                    'x': las.x.astype(np.float32),
                    'y': las.y.astype(np.float32),
                    'z': las.z.astype(np.float32)
                }

                # 如果存在强度信息
                if hasattr(las, 'intensity'):
                    points['intensity'] = las.intensity

                # 计算统计信息
                stats = PointCloudStats(
                    total_points=len(las),
                    bbox={
                        'min': {
                            'x': float(las.header.min[0]),
                            'y': float(las.header.min[1]),
                            'z': float(las.header.min[2])
                        },
                        'max': {
                            'x': float(las.header.max[0]),
                            'y': float(las.header.max[1]),
                            'z': float(las.header.max[2])
                        }
                    },
                    mean_z=float(np.mean(las.z)),
                    max_z=float(np.max(las.z)),
                    min_z=float(np.min(las.z))
                )

                return {
                    'points': points,
                    'stats': stats,
                    'format': 'LAS/LAZ'
                }

        except Exception as e:
            raise Exception(f"解析LAS/LAZ文件失败: {str(e)}")

    @staticmethod
    def parse_xyz(file_path: str) -> Dict[str, Any]:
        """解析XYZ文本格式点云"""

        points = {'x': [], 'y': [], 'z': []}
        bbox = {'min': {'x': float('inf'), 'y': float('inf'), 'z': float('inf')},
                'max': {'x': float('-inf'), 'y': float('-inf'), 'z': float('-inf')}}

        try:
            with open(file_path, 'r') as f:
                for line_num, line in enumerate(f, 1):
                    line = line.strip()
                    if not line or line.startswith('#'):
                        continue

                    parts = line.split()
                    if len(parts) < 3:
                        continue

                    try:
                        x, y, z = map(float, parts[:3])

                        # 更新点数据
                        points['x'].append(x)
                        points['y'].append(y)
                        points['z'].append(z)

                        # 更新包围盒
                        bbox['min']['x'] = min(bbox['min']['x'], x)
                        bbox['min']['y'] = min(bbox['min']['y'], y)
                        bbox['min']['z'] = min(bbox['min']['z'], z)
                        bbox['max']['x'] = max(bbox['max']['x'], x)
                        bbox['max']['y'] = max(bbox['max']['y'], y)
                        bbox['max']['z'] = max(bbox['max']['z'], z)

                    except ValueError:
                        continue

        except Exception as e:
            raise Exception(f"解析XYZ文件失败: {str(e)}")

        # 转换为numpy数组
        points = {
            'x': np.array(points['x'], dtype=np.float32),
            'y': np.array(points['y'], dtype=np.float32),
            'z': np.array(points['z'], dtype=np.float32)
        }

        # 计算统计信息
        stats = PointCloudStats(
            total_points=len(points['x']),
            bbox=bbox,
            mean_z=float(np.mean(points['z'])),
            max_z=float(np.max(points['z'])),
            min_z=float(np.min(points['z']))
        )

        return {
            'points': points,
            'stats': stats,
            'format': 'XYZ'
        }

    @staticmethod
    def downsample_points(points: Dict[str, np.ndarray],
                         target_count: int = 1000000,
                         method: str = 'random') -> Dict[str, np.ndarray]:
        """降采样点云数据"""

        current_count = len(points['x'])

        if current_count <= target_count:
            return points

        if method == 'random':
            # 随机采样
            indices = np.random.choice(
                current_count,
                target_count,
                replace=False
            )
        elif method == 'grid':
            # 网格降采样
            voxel_size = np.cbrt(current_count / target_count)

            # 计算网格索引
            x_indices = (points['x'] // voxel_size).astype(int)
            y_indices = (points['y'] // voxel_size).astype(int)
            z_indices = (points['z'] // voxel_size).astype(int)

            # 去重
            unique_indices = np.unique(
                np.column_stack([x_indices, y_indices, z_indices]),
                axis=0
            )

            if len(unique_indices) > target_count:
                # 如果网格点仍然太多，再随机采样
                indices = np.random.choice(
                    len(unique_indices),
                    target_count,
                    replace=False
                )
                indices = unique_indices[indices]
                return {
                    'x': points['x'][indices],
                    'y': points['y'][indices],
                    'z': points['z'][indices]
                }
            else:
                indices = np.arange(current_count)[
                    np.lexsort((z_indices, y_indices, x_indices))
                ][:target_count]
        else:
            raise ValueError(f"不支持的降采样方法: {method}")

        return {
            'x': points['x'][indices],
            'y': points['y'][indices],
            'z': points['z'][indices]
        }

    @staticmethod
    def export_pointcloud(points: Dict[str, np.ndarray],
                         output_path: str,
                         format: str = 'LAS') -> None:
        """导出点云数据"""

        if format.upper() == 'LAS':
            # 创建LAS文件
            header = laspy.LasHeader(point_format=3, version="1.2")
            header.offsets = [
                np.min(points['x']),
                np.min(points['y']),
                np.min(points['z'])
            ]
            header.scales = [0.01, 0.01, 0.01]

            las = laspy.LasData(header)
            las.x = points['x']
            las.y = points['y']
            las.z = points['z']

            if 'intensity' in points:
                las.intensity = points['intensity']

            las.write(output_path)
        else:
            # 导出为XYZ文本
            with open(output_path, 'w') as f:
                for i in range(len(points['x'])):
                    f.write(f"{points['x'][i]} {points['y'][i]} {points['z'][i]}\n")
\end{lstlisting}

\newpage
\bibliographystyle{gbt7714-numerical}
\bibliography{cas-refs}
\newpage
\appendix

\section{部署与配置管理}

\subsection{环境变量配置}
\begin{table}[pos=!htbp]
    \centering
    \caption{环境变量配置说明}
    \label{tab:env-config}
    \begin{tabular}{@{} LLL @{}}
        \toprule
        \textbf{环境变量}                  & \textbf{说明}     & \textbf{默认值}            \\
        \midrule
        \texttt{FEHALS\_HOST}          & 服务器监听地址         & 0.0.0.0                 \\
        \texttt{FEHALS\_PORT}          & 服务器端口           & 8000                    \\
        \texttt{HELIOS\_PATH}          & HELIOS++可执行文件路径 & /usr/local/bin/helios++ \\
        \texttt{HELIOS\_REPO}          & HELIOS++仓库路径    & /opt/helios++/repo      \\
        \texttt{HELIOS\_ASSETS}        & HELIOS++资源目录    & /opt/helios++/assets    \\
        \texttt{FEHALS\_SIM\_TIMEOUT}  & 仿真超时时间（秒）       & 300                     \\
        \texttt{FEHALS\_CORS\_ORIGINS} & 允许的跨域源          & http://localhost:3000   \\
        \texttt{FEHALS\_UPLOAD\_DIR}   & 文件上传目录          & uploads                 \\
        \texttt{FEHALS\_CACHE\_DIR}    & 缓存目录            & cache                   \\
        \bottomrule
    \end{tabular}
\end{table}

\subsection{Docker部署方案}
\begin{lstlisting}[language=bash,caption=Dockerfile - 容器化部署]
# Dockerfile
FROM python:3.9-slim

WORKDIR /app

# 安装系统依赖
RUN apt-get update && apt-get install -y \
    gcc \
    g++ \
    git \
    wget \
    && rm -rf /var/lib/apt/lists/*

# 安装HELIOS++
RUN wget https://github.com/3dgeo-heidelberg/helios_plusplus/releases/download/v1.0.0/helios++_linux_v1.0.0.tar.gz && \
    tar -xzf helios++_linux_v1.0.0.tar.gz && \
    mv helios++ /usr/local/bin/

# 复制requirements文件并安装Python依赖
COPY requirements.txt .
RUN pip install --no-cache-dir -r requirements.txt

# 复制应用代码
COPY . .

# 创建必要的目录
RUN mkdir -p uploads cache static logs

# 设置环境变量
ENV FEHALS_HOST=0.0.0.0
ENV FEHALS_PORT=8000
ENV HELIOS_PATH=/usr/local/bin/helios++
ENV HELIOS_REPO=/app/helios++/repo
ENV HELIOS_ASSETS=/app/helios++/assets

# 暴露端口
EXPOSE 8000

# 启动命令
CMD ["python", "run.py"]
\end{lstlisting}

\subsection{服务监控与日志}
\begin{lstlisting}[language=bash,caption=监控配置示例]
# 使用Prometheus + Grafana监控
# prometheus.yml
global:
  scrape_interval: 15s

scrape_configs:
  - job_name: 'fehals'
    static_configs:
      - targets: ['fehals-backend:8000']
    metrics_path: '/metrics'
    scrape_interval: 15s
\end{lstlisting}

\section{开发规范与最佳实践}

\subsection{代码规范}
\begin{itemize}
    \item \textbf{Python代码}：遵循PEP 8规范，使用black格式化，flake8检查
    \item \textbf{Vue代码}：使用ESLint + Prettier，采用Composition API
    \item \textbf{Git提交规范}：使用Angular提交格式（type(scope): message）
    \item \textbf{分支管理}：master（主分支）、develop（开发分支）、feature/*（功能分支）
\end{itemize}

\subsection{API设计原则}
\begin{itemize}
    \item \textbf{RESTful}：遵循REST设计规范，使用标准HTTP方法
    \item \textbf{状态码}：使用标准HTTP状态码，如200、201、400、404、500等
    \item \textbf{错误处理}：统一错误格式，包含错误代码和描述信息
    \item \textbf{文档化}：所有API接口必须有Swagger/OpenAPI文档
\end{itemize}

\section{故障排查指南}

\begin{table}[pos=!htbp]
    \centering
    \caption{常见问题及解决方案}
    \label{tab:troubleshooting}
    \begin{tabular}{@{} LL @{}}
        \toprule
        \textbf{问题现象} & \textbf{解决方案}                \\
        \midrule
        HELIOS++无法启动  & 检查HELIOS\_PATH环境变量，确认可执行文件存在 \\
        模型文件上传失败      & 检查文件格式和大小，确认在允许范围内           \\
        仿真任务卡住        & 检查HELIOS++资源目录，查看任务日志        \\
        WebSocket连接断开 & 检查CORS配置，确认代理设置正确            \\
        点云渲染性能差       & 使用降采样，减少同时渲染的点数              \\
        \bottomrule
    \end{tabular}
\end{table}

\end{document}